\documentclass[10pt,twocolumn]{article}

\usepackage[utf8]{inputenc}
\usepackage[T1]{fontenc}
\usepackage{lmodern}
\usepackage[margin=2cm]{geometry}
\usepackage{array}
\usepackage{booktabs}
\usepackage{tabularx}
\usepackage{url}
\usepackage[hidelinks]{hyperref}

\setlength{\columnsep}{0.8cm}

\begin{document}

\begin{titlepage}
  \centering
  {\large Faculty of Engineering\par}
  {\large ECCE Department\par}
  \vspace{1.5cm}
  {\Large EEN 527\par}
  {\Large Embedded System Design\par}
  {\large Section A\par}
  \vspace{1cm}
  {\large Instructor:\par}
  {\large Dr.\ Abdallah Kassem\par}
  \vspace{2cm}
  {\LARGE\bfseries Literature Review\par}
  \vspace{0.8cm}
  {\Large FPGA Based Intrusion Detection and\par}
  {\Large Prevention System for CAN Bus Networks\par}
  \vfill
  {\large Submitted by:\par}
  \vspace{0.3cm}
  {\large Anthony El Chakar \#20231589\par}
  {\large Charbel Estephan \#20231222\par}
  {\large Joe Geagea \#20201000\par}
\end{titlepage}

\section{Introduction}

This review covers existing research on intrusion detection and prevention
systems for the Controller Area Network (CAN) protocol, with a focus on
hardware based and FPGA based implementations. CAN networks power most
embedded systems in vehicles and industrial equipment, but the protocol has
no built in authentication or encryption. This gap leaves it open to
spoofing, denial of service, and replay attacks.

The review draws on IEEE Xplore, ScienceDirect, and arXiv, and it covers work
published between 2012 and 2026. It includes conference papers, journal
articles, and preprints on CAN intrusion detection and prevention, whether
the design runs in software, on a microcontroller, on an FPGA, or on an
application specific circuit. It excludes intrusion detection work aimed at
other automotive networks such as FlexRay or Automotive Ethernet, and it
treats message authentication and encryption schemes only in passing, since
they solve the same problem through a different mechanism than the detection
and prevention systems this review examines.

The rest of this document follows a fixed structure. Section~\ref{sec:background}
lays out the background a reader needs on the CAN protocol itself, the
categories of attack the literature studies, and the built in mechanism
every prevention design in this review relies on to stop a bad frame once it
finds one. Section~\ref{sec:related} surveys the related work in three
groups, rule based and signature based detectors, machine learning based
detectors, and hardware implementations of either kind, then closes with the
gap this project targets. Section~\ref{sec:summary} succinctly restates this
limitation and details how the proposed hardware architecture addresses it.

\section{Background}
\label{sec:background}

CAN organizes traffic as broadcast frames rather than addressed packets. Each
frame carries an arbitration field of 11 or 29 bits that both identifies the
message and sets its priority, since the frame with the lowest identifier
wins the bus when several nodes try to transmit at the same time. A control
field states the data length, up to eight bytes of data follow, and a cyclic
redundancy check, an acknowledge field, and an end of frame field close the
frame out. Every node on the bus receives every frame, and the protocol has
no field for a sender address, a signature, or an encryption key. Robert
Bosch GmbH, who introduced the protocol, built it for a closed vehicle
environment where every node was assumed trustworthy. That assumption no
longer holds once a node can be reached through the On Board Diagnostics
port, an infotainment system, or a cellular modem.

Attacks on CAN fall into a small number of recurring categories. Ileri et
al.\ group them as denial of service, fuzzy, malfunction, replay, and
masquerade attacks~\cite{ileri2025}. A denial of service attack floods the
bus with a high priority identifier so that legitimate frames lose
arbitration. A fuzzy attack injects frames with random identifiers and random
payloads to trigger unexpected behavior. A malfunction or spoofing attack
reuses a legitimate identifier with a forged payload, so a receiving node
accepts data it should not. A replay attack retransmits a previously captured
frame. A masquerade attack silences the legitimate sender of an identifier
and takes its place. Miller and Valasek showed that an attacker can reach the
bus without any physical access at all, disabling brakes, altering the
instrument cluster, and taking other unsafe actions through the cellular,
Bluetooth, and Wi Fi interfaces built into modern telematics
units~\cite{miller2015}. Their work moved the field's threat model away from
an attacker who needs a wrench and a laptop under the dashboard and toward
one who needs only an internet connection, and it is one of the main reasons
intrusion detection research on CAN accelerated afterward rather than
treating physical bus access as an unrealistic assumption.

Seo et al.\ built one of the datasets most later work relies on by recording
normal traffic from a Hyundai YF Sonata and injecting denial of service,
fuzzy, and RPM or gear spoofing frames at fixed intervals~\cite{seo2018}.
Because that same vehicle and that same set of attack types recur across so
many of the papers surveyed below, this review treats it as the default
benchmark and calls out any paper that departs from it. Donadel et al.\
separate attacks by the access level they need, distinguishing frame
injection, which any node with bus access can attempt, from single bit
attacks, which require an attacker who can drive individual bits on the
physical lines~\cite{donadel2025}. Araujo Filho et al.\ state the threat
model that most of the papers in this review adopt. Their attacker can sniff
the bus, learn the identifiers and timing patterns in use, inject new frames,
and modify a frame in transit while updating its checksum~\cite{araujo2021}.
This project adopts the same threat model, since it maps directly onto a
device that taps the physical transmit and receive lines, reads every bit as
it passes, and decides whether to act before the frame ends.

The literature groups defenses into two broad classes. Signature based
systems compare traffic against a set of known attack patterns and flag a
match. They detect known attacks reliably but miss anything outside their
signature set. Anomaly based systems learn what normal traffic looks like and
flag deviations from it. This second class covers most of the work reviewed
here, and it splits further into systems that model message timing and
frequency, systems that check a message against the physical or electrical
fingerprint of the sending node, and systems that extract statistical or
learned features from the identifier, timing, and payload fields and classify
them with a machine learning model. A random forest, the model class this
project builds around, sits in that last group. It reads a set of features
computed from a frame, sends the features down many decision trees trained on
different random subsets of the training data, and takes a majority vote
across the trees, which is what gives it a reputation for resisting the
overfitting that a single deep decision tree suffers from.

CAN also gives every defense a built in tool for stopping an attack once it
is caught. The bit stuffing rule requires a change of level after five
consecutive identical bits, so a node can force an error condition by sending
six dominant bits in a row. Every other node reads this as a protocol
violation, discards the frame in progress, and sends its own error frame, and
repeated violations from one node raise its error counters until it enters
the bus off state and stops transmitting. This mechanism matters for a
prevention system for two reasons beyond simply blocking one bad frame.
First, it needs no cooperation from the rest of the bus, since every CAN
controller already implements bit stuffing and error counting as part of the
protocol, so a single added device can trigger it without touching any other
node's firmware. Second, it is available the instant the offending bits are
seen, which is why the timing budget in Section~\ref{sec:hardware} measures
against the frame's own end rather than against some later software response.
Section~\ref{sec:hardware} returns to this mechanism, since most of the
prevention designs in this review use it.

Two adjacent lines of work fall outside detection and prevention but shape
why they matter. Message authentication schemes add a code to each frame so a
receiver can reject anything it cannot verify, at the cost of extra bits that
CAN's small payload struggles to spare. Kurachi et al.\ built one such scheme,
CaCAN, around a centralized node that authenticates messages on behalf of the
rest of the bus~\cite{kurachi2014}. Physical layer schemes such as bitrate
hopping change the bus speed on a shared secret sequence so an
unsynchronized attacker cannot follow the traffic at all~\cite{lambert2025}.
Both approaches need every node on the bus to cooperate, which is a heavier
requirement than adding one detection device, and this is part of why
detection and prevention systems that tap the bus from a single point remain
an active research direction.

\section{Related Work}
\label{sec:related}

\subsection{Rule Based and Signature Based Detection}
\label{sec:rules}

The earliest CAN specific detectors work from a small number of hand chosen
features rather than a trained model. Song et al.\ showed that the interval
between successive messages with the same identifier stays stable under
normal operation and that fuzzy and denial of service attacks disturb it
enough to raise no false alarms in their tests~\cite{song2016}. This result
matters beyond its own numbers, since it shows that a feature as simple as a
timestamp difference, something a small microcontroller can compute with a
single subtraction, already carries enough signal to separate two of the
four attack categories from normal traffic without any learned model at all.
Lee et al.\ built on this idea by measuring how far each frame's interval and
payload Hamming distance deviate from the averages recorded for its
identifier in normal traffic, then passing those two deviations to a three
tree random forest trained in WEKA to flag spoofing~\cite{lee2023}. Because
both features come from comparing a new frame against a running average kept
per identifier, the detector needs only a small table indexed by identifier
and a handful of arithmetic operations per frame, which is part of why
Section~\ref{sec:hardware} shows this same design translating cleanly into
dedicated hardware.

Detection based on message ownership offers a different rule. Abbott McCune
and Shay gave every node a whitelist of the identifiers it owns and flagged
any frame carrying one of those identifiers while the owning node was not
transmitting, catching a replay directly without needing a learned
model~\cite{abbott2016}. Pese et al.\ built MichiCAN on the same principle at
the bit level, monitoring the arbitration field bit by bit from inside the
CAN controller itself, which lets it flag a denial of service attack before
the identifier field finishes rather than waiting for the whole
frame~\cite{pese2025}. Donadel et al.\ take the most direct rule of all with
CANTXSec, an officer that watches which electronic control unit drives the
transmit line for each identifier and flags any frame whose sender does not
match the expected node, without any learning step at all~\cite{donadel2025}.

These designs share a strength and a weakness. Because they check a fixed
rule rather than a learned pattern, their behavior is easy to reason about
and their detection is immediate once the relevant bits arrive. Because the
rule is fixed, they only catch what their designer anticipated, and Donadel
et al.\ note that a purely rule based officer needs no training data at all,
in contrast with every machine learning design discussed
next~\cite{donadel2025}. This tradeoff is why rule based detection alone does
not fit a project built around a random forest. A whitelist or an ownership
rule catches masquerade and replay attacks cleanly but has nothing to say
about a fuzzy attack that forges a payload for an identifier the correct node
is still transmitting, which is exactly the gap a learned classifier is meant
to close.

\subsection{Machine Learning Based Detection}

Most recent CAN intrusion detectors apply a classifier to features drawn from
the identifier, the timing, or the payload. Seo et al.\ converted CAN
identifiers into a simple image and trained GIDS, a pair of generative
adversarial discriminators, where the first discriminator learns known
attacks and the second learns only from normal traffic so it can flag
attacks it never saw~\cite{seo2018}. Their design shows that even a shallow,
hand engineered transformation of the identifier field carries enough
structure for a neural network to work with, though the two discriminator
training loop that GIDS needs adds cost that a simpler tree based classifier
avoids.

Some of the reviewed designs train only on normal traffic, which lets them
flag attacks without ever seeing a labeled example, while others train on
labeled attacks. Araujo Filho et al.\ compared a one class support vector
machine against an isolation forest and found the isolation forest reached
accuracy above 99 percent and an F1 score above 97 percent while reading only
the first five payload bytes of each frame, which matters because fewer bytes
read means less time spent before a decision can be made~\cite{araujo2021}.
Ileri et al.\ built MetaCAN, an XGBoost classifier tuned with a hybrid of
particle swarm optimization and cuckoo search, and reported accuracies from
0.984 to 0.99999 across three public datasets, though training took between
roughly three and ninety four hours depending on the
dataset~\cite{ileri2025}. XGBoost and random forest both build an ensemble of
decision trees, but XGBoost grows each new tree to correct the errors of the
ones before it, while a random forest trains every tree independently on a
random subset of the data and simply averages their votes. The dependence
between XGBoost trees exists only during training. Once trained, both models
can evaluate every tree in parallel, since a random forest combines its trees
by a vote and XGBoost combines them by a sum. The difference that matters for
hardware is in that last step. Each random forest leaf holds a class label, so
combining the trees needs only a small vote counter, while each XGBoost leaf
holds a real valued score, so combining them needs an adder tree and a final
comparison.

These results rest on a small set of shared benchmarks. The Car Hacking
Dataset, recorded from the same Hyundai Sonata used by Seo et al., remains
the default benchmark for denial of service, fuzzy, and spoofing
attacks~\cite{seo2018}. Because nearly every paper in this section reports
accuracy on one or more of these same datasets, the differences in reported
accuracy trace mainly to whether a model was trained on the specific attack
it was later tested against, not to which algorithm was used, a pattern that
carries directly into the hardware implementations discussed next.

\subsection{Hardware Implementations}
\label{sec:hardware}

A growing body of work moves these classifiers off a general purpose
processor and onto dedicated hardware, chiefly to meet the timing budget that
a CAN frame allows. Araujo Filho et al.\ derive that budget directly from the
protocol. If a detector reads the first $N$ of eight payload bytes, the time
left to decide and to send an error frame before the current frame's end of
frame field arrives equals $19 + 8(8 - N)$ bit times, minus whatever delay
the hardware itself adds. At 500 kilobits per second and five payload bytes
this comes to about 86 microseconds, and their isolation forest running on a
Raspberry Pi 4 averaged 78 to 79 microseconds across the three attacks they
tested, meeting the budget with margin to spare~\cite{araujo2021}. That margin
came from software running on a general purpose processor with an operating
system underneath it, and it sets the baseline every hardware design
discussed below is trying to beat by a wider margin, at a lower power draw,
or both.

Khandelwal and Shreejith built an FPGA based intrusion detection system on a
hybrid device that combines an ARM processor with programmable logic acting
as a single electronic control unit. Their design, a quantized convolutional
network reading a window of four identifiers, reaches an F1 score of 0.9996
for denial of service and 0.9922 for fuzzing at 0.43 milliseconds per frame
while drawing 2.0 watts, against 54 watts for the same model on a
GPU~\cite{khandelwal2022}. The quantization step, shrinking every weight down
to a few bits instead of the 32 bit floating point values a GPU trains with,
is what lets a neural network of this size fit onto programmable logic at
all, and it is the main reason the power gap against the GPU is as large as
it is.

Other groups report similar gains. Rangsikunpum et al.\ paired a binarized
network with a generative adversarial discriminator on a low cost Zedboard,
reaching 0.169 milliseconds per decision at 2.09 watts using no dedicated
multiplier blocks at all~\cite{rangsikunpum2024}. Avoiding dedicated
multiplier blocks matters for a board like the DE10 Lite, since a design that
leans on ordinary programmable logic instead of a fixed number of hardware
multipliers has more headroom to grow without running into a hard resource
ceiling.

Tree based designs follow a separate and generally lighter path. Lee et al.\
synthesized a random forest trained with WEKA directly into Verilog hardware
and reached 98.2 percent accuracy in simulation with three
trees~\cite{lee2023}. Because a decision tree is just a sequence of
comparisons against stored thresholds, each node maps onto a small comparator
circuit rather than the multiply and accumulate operations a neural network
needs, which keeps resource use low even before any quantization step is
applied. A later paper by largely the same group integrates a five tree,
depth nine random forest directly into a CAN controller and fabricates the
result as a 180 nanometer application specific circuit, averaging 1.86
microseconds per detection using 4437 lookup tables against 56733 for a
comparable neural network design on the same device~\cite{lee2025}. That
thirteen fold difference in lookup table count, for two designs solving the
same classification problem on the same fabrication process, is the clearest
evidence in this review that model choice, not just platform choice, drives
resource cost. Jang et al.\ built a lightweight convolutional network sized
for a single CAN frame, cutting parameter count from 4045 to 1789 against
their own earlier design, and report a decision 5.68 microseconds after the
last payload byte arrives, fabricated in a 28 nanometer process at 34.68
nanojoules per inference~\cite{jang2026}.

Abbott McCune and Shay had already split a CAN defense across the same two
kinds of device this project compares. Their prototype runs the ownership
rule from Section~\ref{sec:rules} as VHDL on a small Altera FPGA board that
follows the receive and transmit lines bit by bit, pairs it with an Arduino
Uno microcontroller that makes decisions for the node, and lets the node
force a spoofed frame to error out~\cite{abbott2016}. They verify the design
in ModelSim and do not report latency or power, and because the rule is
fixed rather than learned, the design gives no information about how a
trained classifier behaves on the same split.

Table~\ref{tab:hardware} collects the platform, model, and reported cost for
the hardware designs above. The latency figures are not directly comparable,
since the studies measure from different starting points, some from the
arrival of the first payload byte, others from the last, and one from the end
of the data field, and a reader comparing raw numbers across rows should keep
that difference in mind rather than reading a smaller number as simply faster
hardware.

\begin{table*}[t]
  \centering
  \caption{Reported platform, model, latency, and power or energy across the
  hardware designs reviewed above.}
  \label{tab:hardware}
  \begin{tabularx}{\textwidth}{>{\raggedright\arraybackslash}X
                                >{\raggedright\arraybackslash}X
                                >{\raggedright\arraybackslash}X
                                >{\raggedright\arraybackslash}X
                                >{\raggedright\arraybackslash}X}
    \toprule
    Study & Platform & Model & Latency & Power or Energy \\
    \midrule
    Araujo Filho et al.~\cite{araujo2021} & Raspberry Pi 4 & Isolation forest
      & 78 to 79 microseconds & Not reported \\
    Khandelwal and Shreejith~\cite{khandelwal2022} & Ultra96 V2 & Quantized CNN
      & 0.43 ms per frame & 2.0 W \\
    Rangsikunpum et al.~\cite{rangsikunpum2024} & Zedboard & BNN plus GAN
      & 0.169 ms & 2.09 W \\
    Lee et al.~\cite{lee2025} & FPGA and 180 nm ASIC & Random forest
      & 1.86 microseconds & Not reported \\
    Jang et al.~\cite{jang2026} & 28 nm ASIC & Lightweight CNN
      & 5.68 microseconds after data field & 34.68 nJ per inference \\
    Abbott McCune and Shay~\cite{abbott2016} & Altera FPGA plus Arduino Uno
      & Identifier ownership rule & Not reported & Not reported \\
    \bottomrule
  \end{tabularx}
\end{table*}

Several of the hardware designs above also stop the traffic they flag rather
than only reporting it, using the same error frame mechanism described in
Section~\ref{sec:background}. Araujo Filho et al.\ force six dominant bits
from a Raspberry Pi through an MCP2551 transceiver~\cite{araujo2021}. Lee et
al.\ and Jang et al.\ build the same response into their controller
integrated circuits, and Lee et al.\ show the resulting error flag on an
oscilloscope~\cite{lee2025,jang2026}. Because repeated blocking does not stop
a node from retrying, several of these designs go further and drive the
offending node into the bus off state through repeated error frames rather
than treating one blocked frame as a complete
response~\cite{donadel2025,pese2025,lee2025}. This distinction between
detection and prevention matters for this project's own scope. A design that
only reports an attack still leaves the malicious frame on the bus for every
other node to act on, while a design that forces an error frame in time
removes it before any node can, which is why every build in this project
targets prevention rather than detection alone.

\subsection{Gaps in the Current Literature}

None of the papers reviewed above compares an FPGA only implementation of the
same detection model against a microcontroller only implementation and a
combined FPGA plus microcontroller implementation. The FPGA based designs in
Section~\ref{sec:hardware} report their own latency and power figures in
isolation, most often against a GPU or against each other, and the two
application specific circuit designs integrate their classifier behind an
existing CAN controller rather than testing whether a plain microcontroller
could run the same model within the same timing
budget~\cite{khandelwal2022,lee2025,jang2026}. The rule based, bit level
designs that tap the transmit and receive lines directly, CANTXSec and
MichiCAN, run on a microcontroller rather than an FPGA, and neither paper
tests the same detection logic on both platforms to see where the timing
budget breaks down~\cite{donadel2025,pese2025}. Abbott McCune and Shay pair an
FPGA with a microcontroller, but they use the pair as a single design with a
fixed rule, and they measure neither device on its own~\cite{abbott2016}.
This leaves open a direct, controlled answer to a practical design question.
Given the same random forest model and the same bit level tap into the
physical CAN lines, does an FPGA, a microcontroller, or a combination of the
two give the best balance of detection latency, resource cost, and power
draw. Answering that question needs one model held constant across all three
builds, since comparing an FPGA running one model against a microcontroller
running a different one would leave open whether any measured difference came
from the platform or from the model, which is exactly the confound every
paper surveyed above leaves unresolved.

\section{Summary and Positioning of This Work}
\label{sec:summary}

The literature agrees on two points more than any other. Systems trained on
normal traffic alone, whether an isolation forest, an autoencoder, or a rule
based tap on the transmit line, generalize to attacks they never saw, while
systems trained with labeled attacks reach higher accuracy only on the attacks
they trained on. Once a classifier is committed to an application specific
circuit, detection lands in the single digit microsecond range, well inside
the timing budget the protocol allows. The FPGA designs in
Table~\ref{tab:hardware} took 0.169 and 0.43 milliseconds per decision, both
above the 86 microsecond budget from Section~\ref{sec:hardware}, but both ran
neural networks. So the open question is whether a lighter, tree based model
can bring an FPGA or a microcontroller inside that budget, and which platform
then delivers that speed at the lowest cost and power for a given accuracy
target. The random forest results in Section~\ref{sec:hardware} point toward
an answer, since a tree based model reached microsecond level detection with
an order of magnitude fewer lookup tables than a neural network solving the
same problem, which suggests the model class matters as much to that answer as
the platform does.

This project addresses the gap identified above by building a random forest
classifier, the same model class that two prior hardware papers have already
committed to silicon, and evaluating it on three builds that read the CAN bus
at the bit level through a direct transceiver tap. One build runs entirely on
an FPGA, one runs entirely on a microcontroller, and one splits the work
between the two, with the FPGA handling the time critical tap and decision
and the microcontroller handling logging, thresholds, and the bus off policy.
Measuring detection latency, missed deadlines, accuracy, resource use, and
power across all three builds on the same bus and the same replayed attack
traffic gives the direct, controlled comparison that the reviewed literature
does not provide.

Three open source projects shaped the toolchain for this three way
comparison. The IDS ML repository supplies tree based training code for
vehicle intrusion detection~\cite{idsml}. Its notebooks load the CICIDS2017
dataset, so this project wrote its own loader and training script for the Car
Hacking Dataset. The FPGA random forest project generates Verilog from a
trained scikit learn forest~\cite{snyder}, but it averages one class's leaf
values instead of voting on a class label, and it no longer runs on current
scikit learn releases. This project therefore wrote its own generator, which
produces a pipelined voting classifier that accepts one frame per clock. The
emlearn library converts the same forest into portable C for the
microcontroller~\cite{emlearn}, after a small adjustment to the split
thresholds so that its comparisons match scikit learn exactly. Following
Araujo Filho et al., the model reads only the identifier, the data length,
and the first five payload bytes, which keeps the timing budget at the 86
microseconds derived in Section~\ref{sec:hardware}~\cite{araujo2021}. A
shared set of test frames confirms that the FPGA logic, the C code, and the
Python model return the same class for every frame, so all three builds run
exactly one model.

\begin{thebibliography}{99}

\bibitem{ileri2025}
K.~Ileri, A.~Rakib, and S.~Djahel, MetaCAN, an optimized adaptive hybrid
metaheuristic based intrusion detection system for CAN bus security,
\emph{Vehicular Communications}, vol.~55, article 100956, 2025.

\bibitem{miller2015}
C.~Miller and C.~Valasek, Remote exploitation of an unaltered passenger
vehicle, Black Hat USA, 2015.

\bibitem{seo2018}
E.~Seo, H.~M. Song, and H.~K. Kim, GIDS, GAN based intrusion detection system
for in vehicle network, in \emph{Proceedings of the 16th Annual Conference on
Privacy, Security and Trust}, 2018, pages 1 to 6, doi
10.1109/PST.2018.8514157.

\bibitem{donadel2025}
D.~Donadel, K.~Balasubramanian, A.~Brighente, B.~Ramasubramanian, M.~Conti,
and R.~Poovendran, CANTXSec, a deterministic intrusion detection and
prevention system for CAN bus monitoring ECU activations, arXiv preprint
arXiv 2505.09384, 2025.

\bibitem{araujo2021}
P.~F. de Araujo Filho, A.~J. Pinheiro, G.~Kaddoum, D.~R. Campelo, and
F.~L. Soares, An efficient intrusion prevention system for CAN, hindering
cyber attacks with a low cost platform, \emph{IEEE Access}, vol.~9, pages
166855 to 166869, 2021.

\bibitem{song2016}
H.~M. Song, H.~R. Kim, and H.~K. Kim, Intrusion detection system based on the
analysis of time intervals of CAN messages for in vehicle network, in
\emph{Proceedings of the International Conference on Information Networking},
2016, pages 63 to 68.

\bibitem{lee2023}
D.~Lee, C.~Han, and S.~Lee, Hardware design of intrusion detection system for
automotive CAN bus using random forest, in \emph{Proceedings of the
International Conference on Electronics, Information, and Communication},
2023, doi 10.1109/ICEIC57457.2023.10049883.

\bibitem{abbott2016}
S.~Abbott McCune and L.~A. Shay, Intrusion prevention system of automotive
network CAN bus, in \emph{Proceedings of the IEEE International Carnahan
Conference on Security Technology}, 2016, doi 10.1109/CCST.2016.7815711.

\bibitem{pese2025}
M.~D. Pese, B.~Gozubuyuk, E.~Andrechek, H.~Olufowobi, M.~Hamad, and
K.~G. Shin, MichiCAN, spoofing and denial of service protection using
integrated CAN controllers, in \emph{Proceedings of the 55th Annual IEEE and
IFIP International Conference on Dependable Systems and Networks}, 2025,
pages 443 to 456.

\bibitem{khandelwal2022}
S.~Khandelwal and S.~Shreejith, A lightweight multi attack CAN intrusion
detection system on hybrid FPGAs, in \emph{Proceedings of the International
Conference on Field Programmable Logic and Applications}, 2022, pages 425 to
429.

\bibitem{rangsikunpum2024}
A.~Rangsikunpum, S.~Amiri, and L.~Ost, BIDS, an efficient intrusion detection
system for in vehicle networks using a two stage binarized neural network on
low cost FPGA, \emph{Journal of Systems Architecture}, vol.~156, article
103285, 2024.

\bibitem{lee2025}
J.~Lee, S.~Park, S.~Shin, H.~Im, J.~Lee, and S.~Lee, ASIC design for real
time CAN bus intrusion detection and prevention system using random forest,
\emph{IEEE Access}, vol.~13, pages 129856 to 129869, 2025.

\bibitem{jang2026}
Y.~Jang, H.~Im, J.~Kim, S.~Kim, E.~Kim, and S.~Lee, Hardware software co
optimized lightweight real time CAN intrusion detection and prevention system
for ECUs, \emph{Electronics}, vol.~15, no.~10, article 2108, 2026.

\bibitem{kurachi2014}
R.~Kurachi, Y.~Matsubara, H.~Takada, N.~Adachi, Y.~Miyashita, and
S.~Horihata, CaCAN, centralized authentication system in CAN, in
\emph{Proceedings of the 14th International Conference on Embedded Security in
Cars}, 2014, pages 1 to 9.

\bibitem{lambert2025}
W.~L. Lambert, S.~Ghafoor, and S.~Hollifield, Bitrate hopping, an intrusion
prevention technique to secure the CAN bus based distributed network, in
\emph{Proceedings of the 24th International Symposium on Parallel and
Distributed Computing}, 2025, doi 10.1109/ISPDC67428.2025.00027.

\bibitem{idsml}
Western OC2 Lab, IDS ML, intrusion detection system development using machine
learning algorithms, GitHub repository,
\url{github.com/Western-OC2-Lab/Intrusion-Detection-System-Using-Machine-Learning}.

\bibitem{snyder}
J.~B. Snyder, FPGA random forest, an FPGA implementation of an SKLearn random
forest that generates Verilog directly from a trained model, GitHub
repository, \url{github.com/johnbensnyder/FPGA_random_forest}.

\bibitem{emlearn}
J.~Nordby, emlearn, machine learning inference engine for microcontrollers and
embedded devices, Zenodo, 2019, doi 10.5281/zenodo.2589394,
\url{github.com/emlearn/emlearn}.

\end{thebibliography}

\end{document}
